\documentclass[a4paper,10pt]{article}

\usepackage[utf8]{inputenc}
\usepackage[T1]{fontenc}
\usepackage{geometry}
\usepackage{booktabs}
\usepackage{array}
\usepackage{tabularx}
\usepackage{longtable}
\usepackage{listings}
\usepackage{xcolor}
\usepackage[hidelinks]{hyperref}
\usepackage{amsmath}


\geometry{
    a4paper,
    top=1.8cm,
    bottom=1.8cm,
    left=1.8cm,
    right=1.8cm
}

\setlength{\parindent}{0pt}
\setlength{\parskip}{5pt}
\renewcommand{\arraystretch}{1.18}


\newcommand{\hex}[1]{\texttt{0x#1}}
\newcommand{\code}[1]{\texttt{#1}}

\title{\textbf{DEOS Steering \& Brake Node ICD}\\
\large v1.0 (PDF Revizyonu)}
\author{DEOS -- Dokuz Eylül Otonom Sistemler}
\date{28 Eylül 2026}

\begin{document}
\maketitle

\section{Fren (Brake) Alt Sistemi (0x12)}

Fren sistemi \code{DEOS\_NODE\_BRAKE} (\hex{12}) adresindedir.

\subsection{Gerçek Zamanlı Kontrol Komutları (Real-Time RX)}

\begin{center}
\begin{tabularx}{\textwidth}{>{\raggedright\arraybackslash}p{3.8cm} >{\raggedright\arraybackslash}p{3cm} >{\raggedright\arraybackslash}p{2.2cm} >{\raggedright\arraybackslash}p{2.2cm} X}
\toprule
\textbf{Değişken Adı} & \textbf{Sınıf \& Servis} & \textbf{Komut ID} & \textbf{Ölçek} & \textbf{Fonksiyonel Görev} \\
\midrule
\code{target\_brake\_demand} & COMMAND (0x1) BRAKE (0x03) & 0x01 & 0.01 \% & Operatör pedalından veya üst kontrolcüden ham fren talebi. \\
\code{target\_deceleration} & COMMAND (0x1) BRAKE (0x03) & 0x01 / 0x03 & 0.01 m/s² & Otonom planlayıcıdan gelen kinematik hedef yavaşlama ivmesi. \\
\code{brake\_control\_mode} & COMMAND (0x1) BRAKE (0x03) & 0x04 & Enum & Fren komutunun hangi modda icra edileceğini belirler (Yüzde, İvme, Basınç). \\
\code{target\_park\_brake} & COMMAND (0x1) BRAKE (0x03) & 0x05 & Enum & Elektronik Park Freni (EPB) ve rampada dur-kalk kilit desteği. \\
\code{brake\_prefill\_cmd} & COMMAND (0x1) BRAKE (0x03) & 0x06 & Boolean & Balata boşluğunu diske sıfırlayarak reaksiyon süresini kısaltır. \\
\code{estop\_brake\_override} & SAFETY (0x0) BRAKE (0x03) & 0x03 & Boolean & Freni tam sıkar; Kelly Pin 10'u opto ile toprağa çekip gazı keser. \\
\bottomrule
\end{tabularx}
\end{center}

\subsection{Çapraz Sistem Dinamik Telemetrisi (Emniyet ve Kayma Algılama)}
\begin{itemize}
    \item \textbf{current\_vehicle\_speed} (0x10): Araç durmadan mekanik park kilidinin atmasını donanımsal olarak engeller.
    \item \textbf{traction\_wheel\_slip\_state} (0x10): Kilitlenme sezildiğinde fren aktüatörüne kaba ABS gevşetmesi uygular.
    \item \textbf{chassis\_longitudinal\_accel} (0x03): Fren kontrol döngüsünde ivme kapalı çevrim geri besleme girdisi sağlar.
\end{itemize}

\subsection{Fren Konfigürasyon Parametreleri}
\begin{center}
\begin{tabularx}{\textwidth}{>{\raggedright\arraybackslash}p{3.8cm} >{\raggedright\arraybackslash}p{3cm} >{\raggedright\arraybackslash}p{2.2cm} >{\raggedright\arraybackslash}p{2.2cm} X}
\toprule
\textbf{Parametre} & \textbf{ID} & \textbf{Tip} & \textbf{Ölçek} & \textbf{Açıklama} \\
\midrule
PID\_KP/KI/KD & \hex{0001-03} & uint16\_t & 0.01 & Fren pozisyon/basınç kapalı çevrim kazançları. \\
MIN\_POSITION & \hex{0010} & uint16\_t & 0.01 mm & Balatanın mekanik tam serbest kalma sıfır konumu. \\
MAX\_POSITION & \hex{0011} & uint16\_t & 0.01 mm & Balatanın azami basma strok sınırı. \\
CURRENT\_LIMIT & \hex{0012} & uint16\_t & 0.1 A & Aktüatör motorunun termal koruma akım tavanı. \\
SLEW\_RATE\_LIMIT & \hex{0013} & uint16\_t & 0.01 \%/ms & Mekanik dişli şoklarını önleyen maksimum basma hızı. \\
CMD\_TIMEOUT\_MS & \hex{0014} & uint16\_t & 1 ms & İletişim koptuğunda acil duruş frenine geçme gecikmesi. \\
KISS\_POINT\_OFFSET & \hex{0015} & uint16\_t & 0.01 mm/mV & Balatanın diske temas ettiği boşluksuz sıfır noktası. \\
MAX\_FORCE\_LIMIT & \hex{0016} & uint16\_t & 1N / 0.1bar & Mekanik aşırı yük koruma kuvvet tavanı. \\
\bottomrule
\end{tabularx}
\end{center}

\section{Direksiyon (Steering) Alt Sistemi (0x11)}

Direksiyon sistemi \code{DEOS\_NODE\_STEERING} (\hex{11}) adresindedir.

\subsection{Gerçek Zamanlı Kontrol Komutları (Real-Time RX)}
\begin{center}
\begin{tabularx}{\textwidth}{>{\raggedright\arraybackslash}p{3.8cm} >{\raggedright\arraybackslash}p{3cm} >{\raggedright\arraybackslash}p{2.2cm} >{\raggedright\arraybackslash}p{2.2cm} X}
\toprule
\textbf{Değişken Adı} & \textbf{Sınıf \& Servis} & \textbf{Komut ID} & \textbf{Ölçek} & \textbf{Fonksiyonel Görev} \\
\midrule
\code{target\_steering\_angle} & COMMAND (0x1) STEERING (0x02) & 0x01 & 0.01 deg & Hedef tekerlek açısı (Negatif: Sol, Pozitif: Sağ). \\
\code{target\_angular\_velocity} & COMMAND (0x1) STEERING (0x02) & 0x01 & 0.01 deg/s & Direksiyonun hedef açıya döneceği açısal hız tavanı. \\
\code{steering\_control\_mode} & COMMAND (0x1) STEERING (0x02) & 0x04 & Enum & Direksiyon motorunun kontrol modunu belirler. \\
\code{emergency\_stop\_cmd} & SAFETY (0x0) SYSTEM (0x00) & 0x05 & Boolean & Direksiyon motorunu mekanik serbest bırakma (free-wheel). \\
\bottomrule
\end{tabularx}
\end{center}

\subsection{Çapraz Sistem Dinamik Telemetrisi}
\begin{itemize}
    \item \textbf{current\_vehicle\_speed} (0x10): Hız arttıkça azami direksiyon açısını kısarak yüksek hızda devrilmeyi önler.
    \item \textbf{chassis\_yaw\_rate} (0x03): Araç savrulma anında ESP mantığı ile karşı direksiyon (counter-steering) desteği sağlar.
\end{itemize}

\subsection{Direksiyon Konfigürasyon Parametreleri}
\begin{center}
\begin{tabularx}{\textwidth}{>{\raggedright\arraybackslash}p{3.8cm} >{\raggedright\arraybackslash}p{3cm} >{\raggedright\arraybackslash}p{2.2cm} >{\raggedright\arraybackslash}p{2.2cm} X}
\toprule
\textbf{Parametre} & \textbf{ID} & \textbf{Tip} & \textbf{Ölçek} & \textbf{Açıklama} \\
\midrule
PID\_KP / KI / KD & \hex{0001-03} & uint16\_t & 0.01 & Açı kontrolü PID kazançları. \\
PID\_I\_MAX & \hex{0004} & uint16\_t & 0.01 \% & İntegratör doyum patlamasını (windup) önler. \\
FEEDFORWARD\_K & \hex{0005} & uint16\_t & 0.01 & Hedef açısal hız ileri besleme kazancı. \\
DEADBAND & \hex{0006} & uint16\_t & 0.01 deg & Düz gidişte titreşimleri kesen ölü bölge. \\
ZERO\_OFFSET & \hex{0010} & int16\_t & 0.01 deg & Mekanik düz gidiş merkez açı sapması (Kalibrasyon). \\
MIN/MAX\_ANGLE & \hex{0011-12} & int16\_t & 0.01 deg & Yazılımsal sol ve sağ azami dönüş sınırları. \\
CURRENT\_LIMIT & \hex{0013} & uint16\_t & 0.1 A & Direksiyon motoru elektriksel akım tavanı. \\
MAX\_ANGULAR\_VEL & \hex{0014} & uint16\_t & 0.01 deg/s & İzin verilen azami açısal dönüş hızı tavanı. \\
MAX\_ANGULAR\_ACC & \hex{0015} & uint16\_t & 0.01 deg/s² & Mekanik dişliyi koruyan açısal ivme limiti. \\
SOFT\_STOP\_MARGIN & \hex{0016} & uint16\_t & 0.01 deg & Mekanik dayamaya sert çarpmayı önleyen pay. \\
MAX\_TRACK\_ERR & \hex{0017} & uint16\_t & 0.01 deg & Takip hatası toleransı; aşılırsa hata verilir. \\
TRACK\_TIMEOUT\_MS & \hex{0018} & uint16\_t & 1 ms & İzleme hatasının sürdürülebileceği azami süre. \\
\bottomrule
\end{tabularx}
\end{center}

\section{Kalibrasyon Prosedürleri}
Her iki sistem de \code{DEOS\_SERVICE\_CALIBRATION (0x05)} servisini kullanır.
\begin{itemize}
    \item \textbf{Fren:} Balata temas noktası akım sıçraması okuması ile tespit edilir. Kaliper açıkken dara (Tare) alınır.
    \item \textbf{Direksiyon:} Lazerli terazi ile tam düz konum alınarak enkoder açısı \code{ZERO\_OFFSET} olarak EEPROM'a yazılır. Motor düşük torkla sola ve sağa dayamalara sürülerek nominal akım eşiğinde durduğu noktalar yazılımsal sınır açıları (\code{MIN\_ANGLE}, \code{MAX\_ANGLE}) atanır.
\end{itemize}

\end{document}
